\section{Confirmatory results at a glance}
\label{sec:summary}

Table~\ref{tab:summary} lists the three primary contrasts and the two registered secondary contrasts that have an equivalence reading; Table~\ref{tab:readings} gives, for every registered interpretation row that the results triggered, the claim it supports and the claims that do not follow. Sections~\ref{sec:e1} to~\ref{sec:e3} explain each result. The primary contrasts use the \R{level.family}\% family interval; the secondary contrasts use the individual \R{level.ci95}\% interval. The equivalence interval is the \R{level.tost}\% two one-sided test interval for H1 and the \R{level.thetaTtost}\% interval for the secondary contrasts; for the level 2 versus 11 contrast, which had fewer than twenty puzzles with a nonzero difference, both intervals are the registered rare-event bound.

\begin{table}[ht]
\centering
\caption{Confirmatory and registered secondary contrasts. Brier: difference in Brier loss (H1: permuted control minus generated state, positive when the generated state reads better; level 12: teacher minus generated). pp: difference in solve share, in percentage points. Threshold: practical threshold (H2, H3) or equivalence range.}
\label{tab:summary}
\footnotesize
\adjustbox{max width=\linewidth}{\begin{tabular}{p{3.3cm}lrllr}
\toprule
Contrast & Scale & Estimate & Interval & Equiv.\ interval & Threshold \\
\midrule
H1: level-6 readability & Brier & 0.0029 & [0.0025, 0.0033] & [0.0025, 0.0033] & $\pm$0.01 \\
H2: level 1 vs.\ 11 & pp & \ensuremath{-}11.1 & [\ensuremath{-}13.6, \ensuremath{-}8.8] & -- & 5.0 pp \\
H3: compensation & pp & 1.7 & [0.1, 3.3] & -- & 3.0 pp \\
Level 2 vs.\ 11 (secondary) & pp & 0.0 & [\ensuremath{-}1.5, 1.5] & [\ensuremath{-}1.5, 1.5] & $\pm$3.0 pp \\
Level 12, teacher vs.\ generated (secondary) & Brier & 0.000087 & [\ensuremath{-}0.00013, 0.00029] & [\ensuremath{-}0.000091, 0.00026] & $\pm$0.01 \\
\bottomrule
\end{tabular}
}
\end{table}

\begin{table}[ht]
\centering
\caption{Registered interpretation of every triggered row: the result that triggered it, the claim it supports and the claims that do not follow from it (verbatim from the registered interpretation table).}
\label{tab:readings}
\footnotesize
\adjustbox{max width=\linewidth}{\begin{tabular}{p{1.3cm} p{2.9cm} p{3.2cm} p{3.9cm}}
\toprule
 & Registered result & Justified claim & Does not follow \\
\midrule
H1 & H1: 98.33\% lower limit > 0 and the 96.67\% interval inside $\pm$0.01 & The candidate mask is readable from the level-6 state above the permuted null control, but by less than the practical threshold, with the examined probe family and budget & The representation computes candidates, or carries no candidate information; another readout could not do better \\
H1 & Generated probe beats the null control & The label is readable with the examined readout & The model contains more information than the input, or the generator uses the read-out feature \\
H2 & Early single-code replacement causes greater degradation & Greater level sensitivity under the given policy & The coarse code explicitly plans an algorithm; a level 1 vs. 11 difference implies a graded sensitivity along the early levels or is independent of level 1's single-code size and learned upsampler \\
H3 & D better than C, corrected for sham & Later generation attenuates the intervention effect & The improvement is certainly logical error correction \\
H3 & Positive compensation in a checkpoint trained with level corruption & Robustness consistent with the training objective & The model spontaneously detects and repairs logical errors \\
Level 2 vs.\ 11 & Levels 2 vs. 11 interval inside $\pm$3 pp with adequate precision & No practically relevant solve difference between levels 2 and 11 under the given policy & Level 2 carries no information or has no causal role \\
Level 12 & Level-12 readout clearly above the null control and teacher $-$ generated inside $\pm$0.01 & The level-12 readability is consistent with solution content & The generated representation computes candidates or encodes inference steps \\
\bottomrule
\end{tabular}
}
\end{table}
